\documentclass[11pt]{article}
\usepackage{fontspec}
\usepackage[french]{babel}
\usepackage{amsmath}
\usepackage{amssymb}
\usepackage{amsthm}
\usepackage{mathrsfs}
\usepackage{tikz-cd}
\usepackage{enumitem}

\newcommand{\printedpage}[1]{\marginpar{\footnotesize p.~#1}}
\newcommand{\msp}[1]{\marginpar{\footnotesize [#1]}}
\newcommand{\cocart}{\text{\scriptsize cocart.}}

% footnote marks printed as superscript (n) in the text, as superscript n in the footnote
\makeatletter
\AtBeginDocument{%
  \renewcommand\@makefnmark{\hbox{\textsuperscript{\normalfont(\@thefnmark)}}}%
  \renewcommand\@makefntext[1]{\parindent 1em\noindent\hbox to 1.8em{\hss\textsuperscript{\normalfont\@thefnmark}}#1}%
}
\makeatother

\begin{document}

\printedpage{9}
% head: A. Grothendieck : Dérivateurs, Ch. XVII, Édition des Universités de Montpellier II et Paris VII
% foot: Version 9 février 2019
\[
\begin{tikzcd}
X_J \arrow[r] \arrow[d, "u_J"'] \arrow[dr, phantom, "\cocart"] & X_K \arrow[r] \arrow[d, "\tilde{u}_{J,K}"'] \arrow[dr, phantom, "\cocart"] & X_L \arrow[d, "\tilde{u}_{J,L}"'] \arrow[dd, bend left=40, "\tilde{u}_{K,L}"] \\
Y_J \arrow[r] & Y_{J,K} \arrow[r] \arrow[d, "u_{J,K}"'] \arrow[dr, phantom, "\cocart"] & Y_{J,L} \arrow[d, "u^L_{J,K}"'] \arrow[dd, bend left=40, "u_{J,L}"] \\
 & Y_K \arrow[r] & Y_{K,L} \arrow[d, "u_{K,L}"'] \\
 & & Y_L\ ,
\end{tikzcd}
\tag*{(2)}
\]
donnant lieu à une nouvelle flèche $u^L_{J,K}$, déduite de $u_{J,K}$ par le cochangement de base \msp{page 15} $Y_{J,K} \longrightarrow Y_{J,L}$. Nous utiliserons ce diagramme dans les deux cas où $(J \subset K \subset L)$ est l'un des deux triplets
\[
\left\{
\begin{aligned}
& I'_0 \hookrightarrow I_0 \hookrightarrow I \\
& I'_0 \hookrightarrow I' \hookrightarrow I,
\end{aligned}
\right.
\]
issus du carré cocartésien (dans Cat) (carré MV)
\[
\begin{tikzcd}
I'_0 \arrow[r, hook] \arrow[d, hook] & I_0 \arrow[d, hook] \\
I' \arrow[r, hook] & I\ .
\end{tikzcd}
\tag*{(3)}
\]

Il faut insérer les deux diagrammes correspondants du type $(*)$ \texttt{[de la page 10]} dans un même diagramme. On notera à ce propos que les carrés dans $\mathcal{M}$ déduits de (3) en lui appliquant $J \longmapsto X_J$ et $J \longmapsto Y_J$
\[
\begin{matrix} (4_X) \\[12ex] (4_Y) \end{matrix}
\qquad\qquad
\begin{tikzcd}[column sep=small, row sep=small]
X_{I'_0} \arrow[rr] \arrow[dr] \arrow[dd, dotted, "u_{I'_0}"] & & X_{I_0} \arrow[dr] \arrow[dd, dotted, "u_{I_0}"] & \\
 & X_{I'} \arrow[rr] \arrow[dd, dotted, "u_{I'}"'] & & X_I \arrow[dd, dotted, "u_I"'] \\
Y_{I'_0} \arrow[rr] \arrow[dr] & & Y_{I_0} \arrow[dr] & \\
 & Y_{I'} \arrow[rr] & & Y_I
\end{tikzcd}
\qquad\qquad
\left.\vrule height 9ex depth 9ex width 0pt\right\}\ (4)
\]
sont cocartésien dans $\mathcal{M}$. (Transitivité des $\varinjlim$ \ldots). % [sic] cocartésien -> cocartésiens, agreement with « les carrés »

\msp{page 16}
On a un homomorphisme $(u_{I'_0}, u_{I_0}, u_{I'}, u_I)$ déduit de $u$ du carré $(4_X)$ dans les carré $(4_Y)$, % [sic] les carré -> le carré, singular noun with plural article
et cet homomorphisme lui aussi s'insère dans un diagramme, incluant les deux avatars déjà signalé de (2). Voici le diagramme pertinent : % [sic] signalé -> signalés, agreement with « les deux avatars »
\printedpage{10}
% head: A. Grothendieck : Dérivateurs, Ch. XVII, Édition des Universités de Montpellier II et Paris VII
% foot: Version 9 février 2019
% the footnote mark (9) stands in the printed label of the display; \footnotemark/\footnotetext used because \footnote cannot be set inside a display tag
\[
\begin{tikzcd}[column sep=small, row sep=small]
 & X_{I_0} \arrow[rrr] \arrow[dd] & & & X_I \arrow[dd] \\
X_{I'_0} \arrow[ur] \arrow[rr] \arrow[dd] & & X_{I'} \arrow[urr] \arrow[dd] & & \\
 & Y_{I'_0,I_0} \arrow[rrr] \arrow[dd, "u_{I'_0,I_0}"] & & & Y_{I'_0,I} \arrow[dd, "u^I_{I'_0,I_0}"] \arrow[dddl] \\
Y_{I'_0} \arrow[ur] \arrow[rr] \arrow[dr, dotted] \arrow[dddrr, dotted] & & Y_{I'_0,I'} \arrow[urr] \arrow[ddd] & & \\
 & Y_{I_0} \arrow[rrr] \arrow[dddrrr, dotted] & & & Y_{I_0,I} \arrow[ddd, "u_{I_0,I}"] \\
 & & & Y_{I',I} \arrow[ddr, "u_{I',I}"] & \\
 & & Y_{I'} \arrow[ur] \arrow[drr, dotted] & & \\
 & & & & Y_I\ .
\end{tikzcd}
\tag*{(5~\footnotemark[9])}
\]
\footnotetext[9]{diagramme imbitable, à remplacer par diagramme (6) ci-dessus.}

Le cube supérieur est le cube dans $\mathcal{M}$ à 6 faces cocartésiens déduit des trois flèches de $X_{I'_0}$ dans $X_{I_0}$, $X_{I'}$, $Y_{I'_0}$. % [sic] faces cocartésiens -> faces cocartésiennes, gender agreement
On aura besoin de savoir que la \textit{base} du cube est un carré cocartésien. La partie antérieure du diagramme est l'avatar de (2) correspondant aux inclusions $I'_0 \subset I' \subset I$, la partie derrière l'avatar pour l'inclusion $I'_0 \subset I_0 \subset I$. Donc les deux autres carrés verticaux, en dessous du cube supérieur, sont également cocartésiens, ainsi que le carré oblique en pointillés, lequel n'est autre que la base $(4_Y)$ du diagramme (4). De ces caractères du diagramme (5) résulte aisément (quoique \msp{page 17} j'ai mis une bonne journée de gribouillis en bonheur-la-chance, avant de m'en rendre compte enfin \ldots) que le carré en losange à traits pleins
\[
\begin{tikzcd}[column sep=small, row sep=small]
 & Y_{I'_0,I} \arrow[dl] \arrow[dr] & \\
Y_{I',I} \arrow[dr] & & Y_{I_0,I} \arrow[dl] \\
 & Y_I &
\end{tikzcd}
\]
est également cocartésien. On le voit trivialement sur le diagramme déployé
\printedpage{11}
% head: A. Grothendieck : Dérivateurs, Ch. XVII, Édition des Universités de Montpellier II et Paris VII
% foot: Version 9 février 2019
\[
\begin{tikzcd}[column sep=tiny, row sep=small]
 & & Y_{I'_0} \arrow[dl] \arrow[dr] & & \\
 & Y_{I'_0,I'} \arrow[dl] \arrow[dr] & \cocart & Y_{I'_0,I_0} \arrow[dl] \arrow[dr, "u_{I'_0,I_0}"] & \\
Y_{I'} \arrow[dr] & \cocart & Y_{I'_0,I} \arrow[dl] \arrow[dr, "u^I_{I'_0,I_0}"] & \cocart & Y_{I_0} \arrow[dl] \\
 & Y_{I',I} \arrow[dr, "u_{I',I}"] & \cocart & Y_{I_0,I} \arrow[dl] & \\
 & & {?\simeq Y_I \simeq Y_{I'} \amalg_{Y_{I'_0}} Y_{I_0}\ ,} & &
\end{tikzcd} % ?? uncertain: coproduct sign -- \amalg or \sqcup
\tag*{(6)}
\]
qui est finalement le diagramme pertinent (le diagramme (5) ne servait à rien, c'est un mauvais bout par lequel j'ai pris les choses). On voit sur le diagramme (6) que pour prouver $u_{I',I} \in C$ resp. $TC$, il \msp{page 18} suffit de prouver $u_{I'_0,I_0} \in C$ resp. $TC$. Mais c'est là une assertion relative à l'inclusion d'ensembles ordonnés $I'_0 \subset I_0$ \textit{finis.} On est donc réduit à prouver le lemme suivant, où $I_0$ devient $I$ :

\textbf{Lemme.} \textit{Soit $I$ un ensemble ordonné fini, $u_\bullet : X_\bullet \longrightarrow Y_\bullet$ une flèche dans $\mathcal{M}(I^\circ)$, où $\mathcal{M}$ est une catégorie munie d'un ensemble $C \subset \mathrm{Fl}(\mathcal{M})$ de flèches tel que $C$ soit stable par cochangement de base et par composition, et contienne les isomorphismes. On suppose $u_\bullet \in C_I$, i.e. que pour tout $i \in I$, la flèche $u_{I_{<i},I/i}$ est dans $C$. Alors pour toute partie ouverte $I'$ de $I$, la flèche $u_{I',I}$ est dans $C$}~\footnote[10]{\textbf{N.B.} Bien sûr, $u_{J,K} \in C$ si $J \subset K$ sont deux ouverts de $I$, comme on voit en appliquant le lemme à $K$ au lieu de $I$.}\textit{.}

\textbf{Corollaire.} \textit{Si de plus $W \ \subset\ \mathrm{Fl}(\mathcal{M})$ est un localiseur modérément saturé, tel que $TC \overset{\textup{déf}}{=} W \cap C$ soit stable par cochangement de base (et l'est de toutes façons par composition, et contient les isomorphismes), alors on a}
\[
C_I \cap W_{I^\circ} \ = \ (TC)_I,
\]
\textit{donc si $u_\bullet \in C_I$ (cf. ci-dessus) et les $u_i \in W$, alors $u_\bullet \in (TC)_I$, donc en appliquant ce qui est dit pour $C$ à $TC$, on a $u_{I',I} \in TC$ pour toute} \msp{page 19} \textit{partie ouverte $I'$ de $I$.}

On prouvera en même temps le lemme et son corollaire, en procédant par récurrence sur $\operatorname{card} I = n$. Si $n = 0$, i.e. $I = \emptyset$, les énoncés sont triviaux. Par transitivité, en utilisant les diagrammes du type (2), on est ramené au cas où $I = I' \cup \{i_0\}$, où $i_0 \in I \backslash I'$. On pose encore $I_0 = I/i_0$, $I'_0 = I_0 \cap I'$, et le raisonnement qu'on vient de faire nous montre qu'il suffit de prouver que $u_{I'_0,I_0} \in C$ resp. $\in TC$. Si $I'_0 \neq I$, cela résulte de l'hypothèse de récurrence. % [sic] I'_0 \neq I -> I_0 \neq I, the alternative « Sinon, i_0 est un objet final de I » requires I_0 = I
Sinon, $i_0$ est un objet final de $I$, et $I' = I_{<i_0}$, et alors $u_{I',I} \in C$ est inclus dans la définition de $C_I$. Reste à prouver alors que $u_{I',I} \in TC$ si les $u_i \in W$. On a le diagramme
\printedpage{12}
% head: A. Grothendieck : Dérivateurs, Ch. XVII, Édition des Universités de Montpellier II et Paris VII
% foot: Version 9 février 2019
\[
\begin{tikzcd}
X_{I'} \arrow[r] \arrow[d, "u_{I'}"'] \arrow[dr, phantom, "\cocart"] & X_I \arrow[d, "\tilde{u}_{I',I}"'] \arrow[dd, dotted, bend left=40, "u_I"] \\
Y_{I'} \arrow[r] & Y_{I',I} \arrow[d, "u_{I',I}"'] \\
 & Y_I\ ,
\end{tikzcd}
\]
\msp{page 20} et l'hypothèse de récurrence implique $u_{I'} \in TC$, donc $\tilde{u}_{I',I} \in TC$ par cochangement de base. D'autre part, la flèche $u_I : X_I \longrightarrow Y_I$ s'identifie à $X_{i_0} \longrightarrow Y_{i_0}$ ($i_0$ étant élément final de $I$), elle est donc dans $W$. Donc par saturation, $u_{I',I} \in W$, donc $u_{I',I} \in TC$, q.e.d.

Ainsi, la proposition 2 est prouvée, ainsi que son corollaire 1 -- et on voit que si dans la proposition 2 et son corollaire, on suppose $I$ \textit{fini,} alors il est inutile dans les hypothèses pour $C$, resp. $TC = C \cap W$, de supposer la stabilité par composition \textit{transfinie,} la composition ordinaire suffit. Mais il reste à prouver le corollaire 2, que je viens de prouver dans le cas $I$ \textit{fini,} sans avoir à faire appel à une hypothèse ``transfinie''. Il faut maintenant prouver que c'est encore OK dans le cas localement fini, en supposant $C$ stable par composition transfinie. (Il n'y a plus à revenir sur le cas de $W$, car on sait déjà maintenant que $C_I \cap W_{I^\circ} = (TC)_I$, par le corollaire 1 déjà prouvé.)

\msp{page 21}
La démonstration est essentiellement la même que pour le cas particulier déjà traité $I' = \emptyset$, $I'' = I$, et on aurait pu aussi bien rédiger la démonstration dans ce cadre, à l'aide d'une suite strictement croissante transfinie d'ouverts dans $I$, $I_\alpha$, avec \fbox{$I_0 = I',$} $I_{\alpha_0} = I''$ $(= I)$, et
\[
\begin{aligned}
I_\alpha \ &= \ I_{\alpha-1} \cup I/s_{\alpha-1} && \text{si $\alpha$ non limite, $I_{\alpha-1} \neq I$, en choisissant $s_{\alpha-1} \in I \backslash I_{\alpha-1}$,} \\
I_\alpha \ &= \ \textstyle\bigcup_{\beta<\alpha} I_\beta && \text{si $\alpha$ limite,}
\end{aligned}
\]
et en posant encore
\[
X_\alpha \ = \ X_{I_\alpha}, \qquad Y_\alpha \ = \ Y_{I_\alpha}, \qquad u_\alpha = u_{I_\alpha} \ : \ X_\alpha \longrightarrow Y_\alpha.
\]
Soit $\alpha_0$ tel que $I_{\alpha_0} = I$ (indice terminal de la suite), il faut prouver que dans
\[
\begin{tikzcd}
X_0 \arrow[r] \arrow[d] & X_{\alpha_0} \arrow[d] \\
Y_0 \arrow[r] & Y_{0,\alpha_0} \arrow[d, "u_{0,\alpha_0}"] \\
 & Y_{\alpha_0}\ ,
\end{tikzcd}
\]
\printedpage{13}
% head: A. Grothendieck : Dérivateurs, Ch. XVII, Édition des Universités de Montpellier II et Paris VII
% foot: Version 9 février 2019
\msp{page 22} on a
\[
u_{0,\alpha_0} \ \in \ C.
\]

Compte tenu du complément marginal à la proposition 1 (qui se prouve de la même façon \ldots), on est ramené à prouver que l'homomorphisme des systèmes ordinaux envisagés est dans $C_{I_{\alpha_0+1}}$ (où $I_{\alpha_0+1}$ est l'ensemble des ordinaux $\leq \alpha_0$). Or qu'il en soit ainsi se prouve exactement comme tantôt, à cela près qu'on doit d'abord faire une vérification de $(*)$ p.~1 ($\tilde{u}_\bullet \in C$) pour $i_0 = 0$. Alors le diagramme $(*)$ p.~1 est
\[
\begin{tikzcd}
\emptyset \arrow[r] \arrow[d, "\mathrm{iso}"] & X_0 = X_{I'} \arrow[d, "\mathrm{iso}"', "\mathrm{id}"] \\
\emptyset \arrow[r] & X_0 = X_{I'} \arrow[d, "u_0 = u_{I'}"] \\
 & Y_0 = Y_{I'}\ ,
\end{tikzcd}
\]
et il faut donc vérifier $u_0 = u_{I'} \in C$. Or ceci est contenu dans la proposition 2 elle-même, déjà prouvé, qu'on doit donc utiliser. % [sic] prouvé -> prouvée, agreement with « la proposition 2 »
Le reste se prouve comme avant. Ainsi, le corollaire 2 de la proposition 2 est également prouvé.

\textbf{Commentaires :} Il y a un malaise, car les proposition 1 et proposition 2 (ainsi que leurs corollaires) ont essentiellement la même conclusion, avec des hypothèses sur $I$ qui sont \msp{page 23} différentes : bien ordonné dans un cas, localement fini dans un autre. L'idée naturelle, c'est de donner un énoncé commun, qui coiffe les deux à la fois, avec $I$ artinien, i.e. toute suite décroissante est stationnaire. Ce qu'il faut, c'est une hypothèse sur $I$ d'une part, sur $u_\bullet : X_\bullet \longrightarrow Y_\bullet$ dans $\mathcal{M}(I^\circ)$ d'autre part, qui assure ceci :

\begin{enumerate}
\item[$(*)$] $\forall\ i_0 \in I$, et toute partie ouverte $J'$ de $J = I/i_0$, la flèche $u_{J',J}$ dans le diagramme
\[
\begin{tikzcd}
X_{J'} \arrow[r] \arrow[d, "u_{J'}"'] \arrow[dr, phantom, "\cocart"] & X_J \arrow[d, "\tilde{u}_{J',J}"'] \arrow[dd, bend left=40, "u_J"] \\
Y_{J'} \arrow[r] & Y_{J',J} \arrow[d, "u_{J',J}"'] \\
 & Y_J
\end{tikzcd}
\]
est dans $C$.
\end{enumerate}

\textbf{Question :} est-il exacte qu'il suffit, pour la validité de $(*)$, de supposer % [sic] exacte -> exact, agreement with « il »
\begin{enumerate}
\item[a)] $u_\bullet \in C_I$, i.e. $u_{J',J} \in C$ quand $J = I/i_0$, $J' = J \backslash \{i_0\}$, et
\item[b)] \textit{$I$ localement fini ?}
\end{enumerate}

\end{document}
